\section{Approach 2: AI as Predictor (Generative BO)}
\begin{itemize}
    \item A trained VAE encodes known molecules into a 64-dim coordinate map and generates new candidate molecules by stepping around the best-known ones, then decoding.
    \item I paste the candidates to the LLM; it predicts a LogS number for each (no molecule design, just scoring).
    \item A GP fits the landscape over those predictions, and an acquisition function (EI or UCB) picks which candidate to actually test.
    \item The picked molecule is scored with ESOL (uses budget), recorded, and the loop repeats.
\end{itemize}

\subsection{Pipeline}
Unlike Approach 1, the AI never invents a molecule directly here; it only acts as a \textbf{surrogate} that scores candidates a math engine proposes.

\begin{enumerate}
    \item \textbf{VAE (molecule map):} encoder/decoder trained on $\sim$9,000 real molecules. Nearby coordinates decode to similar molecules, so BO can search continuously and still get valid chemistry back.
    \item \textbf{GP (landscape guesser):} fits the solubility surface over the AI's predictions, with uncertainty.
    \item \textbf{Acquisition function:} EI (Expected Improvement) or UCB (Upper Confidence Bound) picks the next candidate, trading off "looks good" vs.\ "I'm unsure here."
\end{enumerate}

Two decision rules (EI, UCB), each run 3--4 times. The first six runs shared one candidate-generation seed; one independent-seed run per rule was added afterward to test robustness.

\subsection{Results}
\begin{table}[h]
\centering
\begin{tabular}{lccc}
\hline
Rule.Run & Best found & MAE & Pairwise accuracy \\
\hline
EI.1  & +1.219 & 0.83 & 0.876 \\
EI.2  & +1.219 & 0.65 & 0.857 \\
EI.3  & +1.219 & 1.39 & 0.844 \\
EI.4  & +1.219 & 1.03 & 0.857 \\
UCB.1 & +1.219 & 0.48 & 0.867 \\
UCB.3 & +1.219 & 0.60 & 0.819 \\
UCB.4 & +0.208 & 1.32 & 0.657 \\
\hline
\end{tabular}
\caption{Results for Approach 2, both acquisition rules, pooled: best found $1.075 \pm 0.354$, pairwise accuracy $0.825 \pm 0.071$.}
\label{tab:results_a2}
\end{table}

Six of seven runs (including EI.4, on a genuinely independent candidate pool) converge on the same invented molecule at +1.219, evidence of a robustly-reachable peak rather than a shared-seed artifact. UCB.4, run on its own independent pool, is a real failure: all 15 picks scored below the seed (+0.208), best-so-far never moved. The AI's raw numbers are shaky (MAE 0.5--1.4), but its pairwise ordering is usually strong (0.82--0.88), which is what BO actually needs to climb.

\begin{figure}[h]
    \centering
    \begin{subfigure}[b]{0.48\linewidth}
        \centering
        \includegraphics[width=\linewidth]{figures/convergence2_gen.png}
        \caption{Best score so far, per evaluation. Six lines step up to +1.22; UCB.4 stays flat at the seed the whole run.}
        \label{fig:convergence2}
    \end{subfigure}
    \hfill
    \begin{subfigure}[b]{0.48\linewidth}
        \centering
        \includegraphics[width=\linewidth]{figures/topk2_gen.png}
        \caption{Top-3-so-far mean per round. Rises and levels off for six runs; flat for UCB.4.}
        \label{fig:topk2}
    \end{subfigure}

    \vspace{0.5em}

    \begin{subfigure}[b]{0.48\linewidth}
        \centering
        \includegraphics[width=\linewidth]{figures/calibration2_gen.png}
        \caption{Predicted vs.\ true LogS. Loose scatter around the diagonal (raw numbers shaky) but trends the right way.}
        \label{fig:calibration2}
    \end{subfigure}
    \hfill
    \begin{subfigure}[b]{0.48\linewidth}
        \centering
        \includegraphics[width=\linewidth]{figures/pairwise2_gen.png}
        \caption{Pairwise accuracy per run vs.\ coin-flip (0.50). Six runs 0.82--0.88; UCB.4 drops to 0.657.}
        \label{fig:pairwise2}
    \end{subfigure}

    \caption{Comparison across Approach 2 runs (EI vs.\ UCB).}
    \label{fig:results_grid_a2}
\end{figure}
